% Markdown SHA-256 69b1c0d37ac09909fdceb178d08235ed04bb127c49d5e847176e51dcb0aea48e

\documentclass[UTF8,fontset=none,zihao=-4,a4paper]{ctexart}
\usepackage[left=20mm,right=20mm,top=16mm,bottom=18mm,footskip=10mm]{geometry}
\usepackage{fontspec,xcolor,amsmath,amssymb,lmodern,fancyhdr,lastpage}
\usepackage[normalem]{ulem}
\usepackage[unicode,colorlinks=true]{hyperref}
\newif\ifTraceMacFonts
\TraceMacFontstrue
\newcommand{\TraceCheckFont}[1]{\IfFileExists{#1}{}{\TraceMacFontsfalse}}
\TraceCheckFont{/System/Library/Fonts/Supplemental/Arial.ttf}
\TraceCheckFont{/System/Library/Fonts/Supplemental/Arial Bold.ttf}
\TraceCheckFont{/System/Library/Fonts/Supplemental/Arial Italic.ttf}
\TraceCheckFont{/System/Library/Fonts/Supplemental/Arial Bold Italic.ttf}
\TraceCheckFont{/System/Library/Fonts/Monaco.ttf}
\TraceCheckFont{/System/Library/Fonts/STHeiti Light.ttc}
\TraceCheckFont{/System/Library/Fonts/STHeiti Medium.ttc}
% 定义 \TraceForcePortableFonts 可在 macOS 上验证无系统字体时的同一分支。
\ifdefined\TraceForcePortableFonts\TraceMacFontsfalse\fi
\ifTraceMacFonts
  \typeout{Trace font branch: macOS TrueType}
  \setmainfont[
    Path=/System/Library/Fonts/Supplemental/,
    BoldFont={Arial Bold.ttf},
    ItalicFont={Arial Italic.ttf},
    BoldItalicFont={Arial Bold Italic.ttf}
  ]{Arial.ttf}
  \setmonofont[Path=/System/Library/Fonts/,Scale=0.86]{Monaco.ttf}
  \setCJKmainfont[
    Path=/System/Library/Fonts/,
    FontIndex=1,
    BoldFont={STHeiti Medium.ttc},
    BoldFeatures={FontIndex=1},
    AutoFakeSlant=0.16
  ]{STHeiti Light.ttc}
  \setCJKmonofont[Path=/System/Library/Fonts/,FontIndex=1]{STHeiti Light.ttc}
\else
  \typeout{Trace font branch: TeX Live Fandol and Latin Modern}
  \setmainfont[
    BoldFont=lmsans10-bold.otf,
    ItalicFont=lmsans10-oblique.otf,
    BoldItalicFont=lmsans10-boldoblique.otf
  ]{lmsans10-regular.otf}
  \setmonofont[Scale=0.90]{lmmono10-regular.otf}
  \setCJKmainfont[
    BoldFont=FandolHei-Bold.otf,
    AutoFakeSlant=0.16
  ]{FandolHei-Regular.otf}
  \setCJKmonofont{FandolHei-Regular.otf}
\fi
\definecolor{TraceInk}{HTML}{233544}
\definecolor{TraceTitle}{HTML}{173046}
\definecolor{TraceMuted}{HTML}{475A6A}
\definecolor{TraceAccent}{HTML}{09645E}
\definecolor{TraceRule}{HTML}{A5B5BB}
\definecolor{TraceURL}{HTML}{154C69}
\hypersetup{urlcolor=TraceURL,linkcolor=TraceURL}
\setlength{\parindent}{0pt}
\setlength{\parskip}{3pt}
\setlength{\emergencystretch}{1.4em}
\raggedbottom
\pagestyle{fancy}
\fancyhf{}
\renewcommand{\headrulewidth}{0pt}
\fancyfoot[L]{\small\color{TraceMuted} 手册逐项审查}
\fancyfoot[R]{\small\color{TraceMuted} \thepage\ /\ \pageref*{LastPage}}
\newcommand{\TraceDocumentTitle}[1]{%
  \par\noindent{\fontsize{24}{31}\selectfont\bfseries\color{TraceTitle}#1\par}%
  \vspace{3pt}\nopagebreak}
\newcommand{\TraceSection}[1]{%
  \par\addvspace{10pt}\noindent{\Large\bfseries\color{TraceTitle}#1\par}\nopagebreak[4]}
\newcommand{\TraceSubsection}[1]{%
  \par\addvspace{7pt}\noindent{\large\bfseries\color{TraceTitle}#1\par}\nopagebreak[4]}
\newcommand{\TraceTableHeading}[1]{%
  \par\noindent{\small\bfseries\color{TraceMuted}#1\par}\nopagebreak[4]}
\newcommand{\TraceItem}[1]{%
  \par{\color{TraceRule}\hrule height 0.40pt}\nobreak\smallskip%
  \noindent{\bfseries\color{TraceTitle}#1\par}\nopagebreak[4]}
\newcommand{\TraceField}[2]{%
  \par\noindent\textbf{\textcolor{TraceAccent}{#1}}\enspace#2\par}
\newbox\TraceCardBox
\newenvironment{TraceCard}
  {\par\addvspace{7pt}\setbox\TraceCardBox=\vbox\bgroup\hsize=\linewidth}
  {\par\egroup
   % 短条目留在同页。长条目交还 TeX 的正常分页，避免整页高度的盒子。
   \ifdim\dimexpr\ht\TraceCardBox+\dp\TraceCardBox\relax<0.30\textheight
     \box\TraceCardBox
   \else
     \unvbox\TraceCardBox
   \fi}
\newcommand{\TraceCode}[1]{\texttt{\mdseries\upshape #1}}
\newenvironment{TraceCodeBlock}
  {\par\smallskip\begingroup\small\ttfamily\raggedright\parskip=0pt}
  {\par\endgroup\smallskip}
\newcommand{\TraceCodeLine}[1]{\noindent\hangindent=1.25em\hangafter=1\strut #1\par}
\newcommand{\TraceLineBreak}{\leavevmode\newline}

\begin{document}
\color{TraceInk}
\hypersetup{pdftitle={AD4080 数字外围逐模块手册审查}}
\hypertarget{trace-anchor-1}{}\pdfbookmark[0]{AD4080 数字外围逐模块手册审查}{trace-heading-1}\TraceDocumentTitle{AD4080 数字外围逐模块手册审查}

这份审查回答两个问题：AD4080 对外承诺了什么，我们的哪段代码确实实现了它。基线为 \href{https://github.com/defineiocc02/20bit\_SAR\_ADC\_Behaviour\_Verification/commit/c912fd8a3796f0e6df4e163070a0797e5355a4e1}{外围提交 c912fd8}。芯片依据为 ADI 的《AD4080 Data Sheet》Rev. B，共 95 页，页码均为印刷页。核对的 PDF 的 SHA-256 为 \TraceCode{f08f2586dc82\penalty150\relax{}d5f583b42071\penalty150\relax{}b3570c06f1fe\penalty150\relax{}736d3053e15a\penalty150\relax{}559a47ed2841\penalty150\relax{}cbc1}。每条手册链接均可直接打开对应页。\href{https://www.analog.com/media/en/technical-documentation/data-sheets/ad4080.pdf}{ADI 原件}\par

\textbf{最重要的差异在系统校正。} 图 96 的信号顺序为 ADC 码 \ensuremath{\rightarrow} 加 OFFSET \ensuremath{\rightarrow} 乘 GAIN \ensuremath{\rightarrow} 滤波选择 \ensuremath{\rightarrow} FIFO 或事件检测。冻结 DSP 先对原码作增益修正，最后加 OFFSET。两种次序在单位增益或零偏移时给出相同结果。将原码置零、OFFSET 写为 +1024、工程增益设为 129/128，两个整数输出分别应为 1032 和 1024。这个输入避开了饱和与半步舍入。\href{https://www.analog.com/media/en/technical-documentation/data-sheets/ad4080.pdf\#page=64}{原图 96} · \href{https://github.com/defineiocc02/20bit\_SAR\_ADC\_Behaviour\_Verification/blob/feat/ad4080-peripherals/subsys/ad4080\_peripherals/calibration/rtl\_verification.json}{逐拍回执}\par

\textbf{怎样读裁定。} “可认可”限于表格明列并已有独立观察的行为。“部分认可”表示同一模块还有明确的产品接线、接口时序或数值缺口。“不认可”表示已有反例，或所宣称的器件功能在产品入口不可用。“工程自定”表示有工程用途，但手册没有指定该内部结构、算法或工具。后两种均不得用来声称复刻了 ADI 内部实现。已知与手册不符的项目另列于末尾。\par

产品实例清单合并三个文件清单，得到 26 个设计模块的声明。另查两块单独验证的保护电路和原 DSP 目录中两个未被产品清单采用的模块。合计 30 个 RTL 模块，逐行列出。主机校准另有 3 个程序。33 行中，3 行的限定功能可认可，19 行部分认可，2 行有明确缺口而不认可，9 行属工程自定。外围目录 34 份 SV 源中的 25 个台架、性质模块及 FPGA 原语模型另列；数字转换组和模拟宏的责任也在末尾交代。\href{https://github.com/defineiocc02/20bit\_SAR\_ADC\_Behaviour\_Verification/blob/feat/ad4080-peripherals/subsys/ad4080\_peripherals/integration/files.f}{服务清单} · \href{https://github.com/defineiocc02/20bit\_SAR\_ADC\_Behaviour\_Verification/blob/feat/ad4080-peripherals/subsys/ad4080\_peripherals/integration/readback\_files.f}{记录清单} · \href{https://github.com/defineiocc02/20bit\_SAR\_ADC\_Behaviour\_Verification/blob/feat/ad4080-peripherals/subsys/ad4080\_peripherals/integration/product/rtl.f}{产品清单} · \href{https://github.com/defineiocc02/20bit\_SAR\_ADC\_Behaviour\_Verification/blob/feat/ad4080-peripherals/subsys/ad4080\_peripherals/interfaces/function\_ownership.csv}{职责清单}\par


\hypertarget{trace-anchor-2}{}\pdfbookmark[0]{1. 系统校准与产品数字输入}{trace-heading-2}\TraceSection{1. 系统校准与产品数字输入}

\TraceTableHeading{模块或条目}
\begin{TraceCard}
\TraceItem{\TraceCode{system\_\penalty150\relax{}cal.\penalty150\relax{}py} \href{https://github.com/defineiocc02/20bit\_SAR\_ADC\_Behaviour\_Verification/blob/feat/ad4080-peripherals/subsys/ad4080\_peripherals/calibration/system\_cal.py}{源码}}
\TraceField{手册规定}{\href{https://www.analog.com/media/en/technical-documentation/data-sheets/ad4080.pdf\#page=71}{p.71 式 (2) 与 System Error Correction Coefficients\penalty150\relax{}} 用两个已知电平求斜率和零点；\href{https://www.analog.com/media/en/technical-documentation/data-sheets/ad4080.pdf\#page=64}{p.64 图 96} 把 OFFSET 加在增益之前。}
\TraceField{代码、实测与尚缺的证据}{两点解斜率有手册依据。已知正弦的峰值和直流中心由片外 ADCToolbox 拟合，这是我们的另一种测量方法。本次修正实物增益表的 OFFSET 起始估计为 −b；显式工程 K/SHIFT 分支按原 DSP 取 −g·b。实物另用两个偏移设置测所选档的灵敏度。\href{https://github.com/defineiocc02/20bit\_SAR\_ADC\_Behaviour\_Verification/blob/feat/ad4080-peripherals/subsys/ad4080\_peripherals/calibration/README.md}{方案和实验}}
\TraceField{裁定}{部分认可}
\end{TraceCard}
\begin{TraceCard}
\TraceItem{\TraceCode{characterize\penalty150\relax{}\_\penalty150\relax{}registers.\penalty150\relax{}py} \href{https://github.com/defineiocc02/20bit\_SAR\_ADC\_Behaviour\_Verification/blob/feat/ad4080-peripherals/subsys/ad4080\_peripherals/calibration/characterize\_registers.py}{源码}}
\TraceField{手册规定}{\href{https://www.analog.com/media/en/technical-documentation/data-sheets/ad4080.pdf\#page=92}{p.92 图 138、表 60} 给出有符号 12 位 OFFSET 和 1 个 ADC LSB 的刻度；\href{https://www.analog.com/media/en/technical-documentation/data-sheets/ad4080.pdf\#page=93}{p.93 图 139、表 61} 给出单位增益及两个端点。}
\TraceField{代码、实测与尚缺的证据}{扫多档 GAIN，用两点斜率比得到有效码表；在所选增益测两次 OFFSET，估计每寄存器步的输出变化。手册支持寄存器和可测量量，扫表策略、CSV 与误差统计为项目方法。尚无本器件的真实扫表数据。\href{https://github.com/defineiocc02/20bit\_SAR\_ADC\_Behaviour\_Verification/blob/feat/ad4080-peripherals/subsys/ad4080\_peripherals/calibration/test\_characterize\_registers.py}{测量格式与用例}}
\TraceField{裁定}{工程自定}
\end{TraceCard}
\begin{TraceCard}
\TraceItem{\TraceCode{verify\_\penalty150\relax{}rtl.\penalty150\relax{}py} \href{https://github.com/defineiocc02/20bit\_SAR\_ADC\_Behaviour\_Verification/blob/feat/ad4080-peripherals/subsys/ad4080\_peripherals/calibration/verify\_rtl.py}{源码}}
\TraceField{手册规定}{\href{https://www.analog.com/media/en/technical-documentation/data-sheets/ad4080.pdf\#page=20}{p.20 图 33、表 7} 给出差分输入与有符号码；\href{https://www.analog.com/media/en/technical-documentation/data-sheets/ad4080.pdf\#page=64}{p.64 图 96} 给出两项校正的次序。}
\TraceField{代码、实测与尚缺的证据}{8,234 个有效输出的对照对象是冻结整数模型。单独加入图 96 判别向量，比对线上 RTL 结果、工程模型和原图推导值；严格模式遇该差异退出 2，普通工程回归单独报告 OPEN\_\penalty150\relax{}MISMATCH。1,024 笔留出实验只测合成通道的系统误差。\href{https://github.com/defineiocc02/20bit\_SAR\_ADC\_Behaviour\_Verification/blob/feat/ad4080-peripherals/subsys/ad4080\_peripherals/calibration/rtl\_verification.json}{逐码结果}}
\TraceField{裁定}{工程自定}
\end{TraceCard}
\begin{TraceCard}
\TraceItem{\TraceCode{sys\_\penalty150\relax{}error\_\penalty150\relax{}corr} \href{https://github.com/defineiocc02/20bit\_SAR\_ADC\_Behaviour\_Verification/blob/7408edf8db9a41eebdd756378b69d2775f107dec/subsys/ad4080\_dsp/rtl/sys\_error\_corr.sv\#L23}{源码}}
\TraceField{手册规定}{\href{https://www.analog.com/media/en/technical-documentation/data-sheets/ad4080.pdf\#page=64}{p.64 图 96} 指定先加后乘；\href{https://www.analog.com/media/en/technical-documentation/data-sheets/ad4080.pdf\#page=71}{p.71 式 (2)} 写出输入输出的仿射关系；\href{https://www.analog.com/media/en/technical-documentation/data-sheets/ad4080.pdf\#page=93}{p.93 表 61} 以 0x200 为单位增益。}
\TraceField{代码、实测与尚缺的证据}{原实现为 x + O + round(x·(G−512)·K/2\textasciicircum{}S)，图 96 对应 g·(x+O)。两种连续模型的差为 (g−1)·O。取 x=0、O=1024、G=0x300、K=1、S=15，RTL 实出 1024，按图应为 1032。两拍流水、舍入和饱和只按项目实现核验。\href{https://github.com/defineiocc02/20bit\_SAR\_ADC\_Behaviour\_Verification/blob/feat/ad4080-peripherals/subsys/ad4080\_peripherals/calibration/rtl\_verification.json}{独立读回}}
\TraceField{裁定}{不认可}
\end{TraceCard}
\begin{TraceCard}
\TraceItem{\TraceCode{ad4080\_\penalty150\relax{}dsp\_\penalty150\relax{}chain} \href{https://github.com/defineiocc02/20bit\_SAR\_ADC\_Behaviour\_Verification/blob/7408edf8db9a41eebdd756378b69d2775f107dec/subsys/ad4080\_dsp/rtl/ad4080\_dsp\_chain.sv\#L22}{源码}}
\TraceField{手册规定}{\href{https://www.analog.com/media/en/technical-documentation/data-sheets/ad4080.pdf\#page=33}{p.33 图 51、52} 与 \href{https://www.analog.com/media/en/technical-documentation/data-sheets/ad4080.pdf\#page=64}{p.64 图 96} 给出校正与旁路、三条滤波分支、事件及记录的相对位置。}
\TraceField{代码、实测与尚缺的证据}{旁路选的是已校正码，三种滤波均以同一校正输出为输入。分流位置符合图示，校正块内部次序有上述反例。补偿路径依赖外部系数。\href{https://github.com/defineiocc02/20bit\_SAR\_ADC\_Behaviour\_Verification/blob/feat/ad4080-peripherals/subsys/ad4080\_peripherals/interfaces/dsp\_review/README.md}{DSP 数据面验证}}
\TraceField{裁定}{部分认可}
\end{TraceCard}
\begin{TraceCard}
\TraceItem{\TraceCode{sar20\_\penalty150\relax{}ad4080\_\penalty150\relax{}top} \href{https://github.com/defineiocc02/20bit\_SAR\_ADC\_Behaviour\_Verification/blob/7408edf8db9a41eebdd756378b69d2775f107dec/subsys/ad4080\_dsp/rtl/sar20\_ad4080\_top.sv\#L23}{源码}}
\TraceField{手册规定}{\href{https://www.analog.com/media/en/technical-documentation/data-sheets/ad4080.pdf\#page=20}{p.20 图 33、表 7} 给出对外转换码格式。}
\TraceField{代码、实测与尚缺的证据}{该适配器可将本工程 offset-\penalty150\relax{}binary 转成有符号补码；样本 ID、故障锁存、停流与恢复是内部交接约定。它是原 DSP 的独立壳，外围产品实际从 20 位补码边界直接驱动 ad4080\_\penalty150\relax{}top。\href{https://github.com/defineiocc02/20bit\_SAR\_ADC\_Behaviour\_Verification/blob/feat/ad4080-peripherals/subsys/ad4080\_peripherals/integration/product/README.md}{产品绑定}}
\TraceField{裁定}{工程自定}
\end{TraceCard}


\hypertarget{trace-anchor-3}{}\pdfbookmark[0]{2. 滤波、同步、事件与 GPIO}{trace-heading-3}\TraceSection{2. 滤波、同步、事件与 GPIO}

\TraceTableHeading{模块或条目}
\begin{TraceCard}
\TraceItem{\TraceCode{sinc1\_\penalty150\relax{}decim} \href{https://github.com/defineiocc02/20bit\_SAR\_ADC\_Behaviour\_Verification/blob/7408edf8db9a41eebdd756378b69d2775f107dec/subsys/ad4080\_dsp/rtl/sinc1\_decim.sv\#L23}{源码}}
\TraceField{手册规定}{\href{https://www.analog.com/media/en/technical-documentation/data-sheets/ad4080.pdf\#page=64}{p.64 图 96、表 28} 给出 sinc1 选择和带宽；\href{https://www.analog.com/media/en/technical-documentation/data-sheets/ad4080.pdf\#page=94}{p.94 表 62，FILTER\_\penalty150\relax{}CONFIG} 给出 N=2 至 1024 及越界钳位。}
\TraceField{代码、实测与尚缺的证据}{逐档对照验证每 N 个有效输入形成一个块均值、输出间隔 N、无效档钳至 1024。认可这项数值抽取。独立的引脚就绪时序和全部 −3 dB 带宽没有在这项检查中测量。\href{https://github.com/defineiocc02/20bit\_SAR\_ADC\_Behaviour\_Verification/blob/feat/ad4080-peripherals/subsys/ad4080\_peripherals/interfaces/dsp\_review/README.md}{逐档结果}}
\TraceField{裁定}{可认可}
\end{TraceCard}
\begin{TraceCard}
\TraceItem{\TraceCode{sinc5\_\penalty150\relax{}decim} \href{https://github.com/defineiocc02/20bit\_SAR\_ADC\_Behaviour\_Verification/blob/7408edf8db9a41eebdd756378b69d2775f107dec/subsys/ad4080\_dsp/rtl/sinc5\_decim.sv\#L23}{源码}}
\TraceField{手册规定}{\href{https://www.analog.com/media/en/technical-documentation/data-sheets/ad4080.pdf\#page=66}{p.66 表 29} 规定五阶滤波的抽取档；\href{https://www.analog.com/media/en/technical-documentation/data-sheets/ad4080.pdf\#page=69}{p.69 Digital Filter Conversion Pulses} 说明建立时间，\href{https://www.analog.com/media/en/technical-documentation/data-sheets/ad4080.pdf\#page=70}{p.70 图 104} 给出频响。}
\TraceField{代码、实测与尚缺的证据}{数字模型按五级积分/梳状滤波并抽取，N=4 的数值输出间隔为 4。测试中第 12 个逻辑 CNV 就出现启动码 0，第 44 个才达到该台架稳态；手册该档建立计数为 26。三者观察对象不同，尚缺按真实 CNV 对齐的首个可用码。\href{https://github.com/defineiocc02/20bit\_SAR\_ADC\_Behaviour\_Verification/blob/feat/ad4080-peripherals/subsys/ad4080\_peripherals/interfaces/dsp\_review/README.md}{测量起点与序列}}
\TraceField{裁定}{部分认可}
\end{TraceCard}
\begin{TraceCard}
\TraceItem{\TraceCode{sinc5\_\penalty150\relax{}comp} \href{https://github.com/defineiocc02/20bit\_SAR\_ADC\_Behaviour\_Verification/blob/7408edf8db9a41eebdd756378b69d2775f107dec/subsys/ad4080\_dsp/rtl/sinc5\_comp.sv\#L23}{源码}}
\TraceField{手册规定}{\href{https://www.analog.com/media/en/technical-documentation/data-sheets/ad4080.pdf\#page=64}{p.64 图 96、表 28} 与 \href{https://www.analog.com/media/en/technical-documentation/data-sheets/ad4080.pdf\#page=94}{p.94 表 62} 规定 sinc5 后再经补偿并作额外 2 倍抽取；\href{https://www.analog.com/media/en/technical-documentation/data-sheets/ad4080.pdf\#page=70}{p.70 图 105–107} 给出响应目标。}
\TraceField{代码、实测与尚缺的证据}{明确绑定工程系数时测过 2N 输出间隔。原厂 FIR 抽头、系数量化、内部舍入无公开值。本产品默认没有补偿系数，不能据工程系数宣称达到 ±0.1 dB 通带纹波。\href{https://github.com/defineiocc02/20bit\_SAR\_ADC\_Behaviour\_Verification/blob/feat/ad4080-peripherals/subsys/ad4080\_peripherals/interfaces/dsp\_review/README.md}{系数和倍率用例}}
\TraceField{裁定}{部分认可}
\end{TraceCard}
\begin{TraceCard}
\TraceItem{\TraceCode{ad4080\_\penalty150\relax{}record\_\penalty150\relax{}guard} \href{https://github.com/defineiocc02/20bit\_SAR\_ADC\_Behaviour\_Verification/blob/feat/ad4080-peripherals/subsys/ad4080\_peripherals/interfaces/dsp\_review/ad4080\_record\_guard.sv\#L24}{源码}}
\TraceField{手册规定}{\href{https://www.analog.com/media/en/technical-documentation/data-sheets/ad4080.pdf\#page=64}{p.64 图 96 与 Digital Filter} 要求换滤波后经 GPIO 同步；\href{https://www.analog.com/media/en/technical-documentation/data-sheets/ad4080.pdf\#page=67}{p.67 图 98}、\href{https://www.analog.com/media/en/technical-documentation/data-sheets/ad4080.pdf\#page=68}{p.68 图 99} 给出同步、就绪及 FIFO 写入关系。}
\TraceField{代码、实测与尚缺的证据}{独立保护电路延后布防，SINC5、N=4、W=1 台架将首个入存码从启动码 0 改为稳态 −1104。实例位于 tb\_\penalty150\relax{}dsp\_\penalty150\relax{}path；产品清单没有这个保护器，真实 0x29 和 0x1C 的直接访问可绕过。\href{https://github.com/defineiocc02/20bit\_SAR\_ADC\_Behaviour\_Verification/blob/feat/ad4080-peripherals/subsys/ad4080\_peripherals/interfaces/dsp\_review/tb\_dsp\_path.sv}{保护用例}}
\TraceField{裁定}{工程自定}
\end{TraceCard}
\begin{TraceCard}
\TraceItem{\TraceCode{event\_\penalty150\relax{}detect} \href{https://github.com/defineiocc02/20bit\_SAR\_ADC\_Behaviour\_Verification/blob/7408edf8db9a41eebdd756378b69d2775f107dec/subsys/ad4080\_dsp/rtl/event\_detect.sv\#L23}{源码}}
\TraceField{手册规定}{\href{https://www.analog.com/media/en/technical-documentation/data-sheets/ad4080.pdf\#page=54}{p.54 图 82、表 26} 与 \href{https://www.analog.com/media/en/technical-documentation/data-sheets/ad4080.pdf\#page=55}{p.55 图 83} 定义内部/外部事件、阈值、ALERT 与状态；\href{https://www.analog.com/media/en/technical-documentation/data-sheets/ad4080.pdf\#page=56}{p.56 图 84、85} 画出迟滞及路由。}
\TraceField{代码、实测与尚缺的证据}{实现上/下阈、共享迟滞；独立数据面测试有低于阈值未触发与高于阈值触发。GPIO 电平、越界脉冲与 0x14 粘滞状态须分别对齐。比较发生在数值 valid，产品尚不能证明从首个越阈 CNV 到外部 ALERT 的两周期约束。\href{https://github.com/defineiocc02/20bit\_SAR\_ADC\_Behaviour\_Verification/blob/feat/ad4080-peripherals/subsys/ad4080\_peripherals/integration/product/evidence/verification.txt}{事件到记录实测}}
\TraceField{裁定}{部分认可}
\end{TraceCard}
\begin{TraceCard}
\TraceItem{\TraceCode{gpio\_\penalty150\relax{}ctrl} \href{https://github.com/defineiocc02/20bit\_SAR\_ADC\_Behaviour\_Verification/blob/7408edf8db9a41eebdd756378b69d2775f107dec/subsys/ad4080\_dsp/rtl/gpio\_ctrl.sv\#L23}{源码}}
\TraceField{手册规定}{\href{https://www.analog.com/media/en/technical-documentation/data-sheets/ad4080.pdf\#page=53}{p.53 表 25} 与 \href{https://www.analog.com/media/en/technical-documentation/data-sheets/ad4080.pdf\#page=87}{p.87 表 52}、\href{https://www.analog.com/media/en/technical-documentation/data-sheets/ad4080.pdf\#page=88}{p.88 表 53}、\href{https://www.analog.com/media/en/technical-documentation/data-sheets/ad4080.pdf\#page=89}{p.89 表 54} 列 4 个 GPIO 的选择、方向、输入/输出含义。}
\TraceField{代码、实测与尚缺的证据}{已有 SDO、FULL、READ\_\penalty150\relax{}DONE、HI/LO、低有效 ALERT、用户 GPO、低有效 SYNC 及外部事件通路。0x3 的 FILTER\_\penalty150\relax{}RESULT\_\penalty150\relax{}RDY 没有接通；保留选择 0xA 在冻结 bank 中可写入。独立控制桥能拒绝该码，产品未实例化桥。\href{https://github.com/defineiocc02/20bit\_SAR\_ADC\_Behaviour\_Verification/blob/feat/ad4080-peripherals/subsys/ad4080\_peripherals/interfaces/control\_review/README.md}{GPIO 用例与边界}}
\TraceField{裁定}{部分认可}
\end{TraceCard}


\hypertarget{trace-anchor-4}{}\pdfbookmark[0]{3. 配置接口、供电命令与数字顶层}{trace-heading-4}\TraceSection{3. 配置接口、供电命令与数字顶层}

\TraceTableHeading{模块或条目}
\begin{TraceCard}
\TraceItem{\TraceCode{spi\_\penalty150\relax{}slave} \href{https://github.com/defineiocc02/20bit\_SAR\_ADC\_Behaviour\_Verification/blob/7408edf8db9a41eebdd756378b69d2775f107dec/subsys/ad4080\_dsp/rtl/spi\_slave.sv\#L24}{源码}}
\TraceField{手册规定}{\href{https://www.analog.com/media/en/technical-documentation/data-sheets/ad4080.pdf\#page=34}{p.34 图 53、54 与表 14}、\href{https://www.analog.com/media/en/technical-documentation/data-sheets/ad4080.pdf\#page=35}{p.35 Write Access 章节} 规定配置帧和普通写数据相 SDO=0；\href{https://www.analog.com/media/en/technical-documentation/data-sheets/ad4080.pdf\#page=41}{p.41 表 18、19} 规定配置 CRC-8。}
\TraceField{代码、实测与尚缺的证据}{长指令、基本顺序流、严格实体和已绑定的 CRC 路径有数字用例。基线产品把普通写相的 SDO 配成高阻，本轮改正并在实际产品引脚断言 OE 和低电平。原件短指令用 7 位地址，产品未传该参数；流式 CRC 的后续实体条件也未全通。\href{https://github.com/defineiocc02/20bit\_SAR\_ADC\_Behaviour\_Verification/blob/feat/ad4080-peripherals/subsys/ad4080\_peripherals/integration/product/evidence/verification.txt}{产品回执} · \href{https://github.com/defineiocc02/20bit\_SAR\_ADC\_Behaviour\_Verification/blob/feat/ad4080-peripherals/subsys/ad4080\_peripherals/interfaces/control\_review/README.md}{控制专项}}
\TraceField{裁定}{部分认可}
\end{TraceCard}
\begin{TraceCard}
\TraceItem{\TraceCode{ad4080\_\penalty150\relax{}regs} \href{https://github.com/defineiocc02/20bit\_SAR\_ADC\_Behaviour\_Verification/blob/7408edf8db9a41eebdd756378b69d2775f107dec/subsys/ad4080\_dsp/rtl/ad4080\_regs.sv\#L23}{源码}}
\TraceField{手册规定}{\href{https://www.analog.com/media/en/technical-documentation/data-sheets/ad4080.pdf\#page=73}{p.73 表 31} 与 \href{https://www.analog.com/media/en/technical-documentation/data-sheets/ad4080.pdf\#page=74}{p.74 表 31 续} 给出地址分布，\href{https://www.analog.com/media/en/technical-documentation/data-sheets/ad4080.pdf\#page=83}{p.83 表 47} 的 0x14 含 POR、ROM 错误及事件状态。}
\TraceField{代码、实测与尚缺的证据}{真实 SPI 可写读 0x25–28 校正系数、0x29 滤波、阈值与水位。0x15–17 只返回默认值并拒绝实际改配置，0x04 产品 ID 低字节在装配中未绑定，0x14 的模拟 POR/ROM 事件未从服务接入。主要寄存器的分项裁定见下文。\href{https://github.com/defineiocc02/20bit\_SAR\_ADC\_Behaviour\_Verification/blob/feat/ad4080-peripherals/subsys/ad4080\_peripherals/interfaces/control\_review/validation.txt}{控制面证据}}
\TraceField{裁定}{部分认可}
\end{TraceCard}
\begin{TraceCard}
\TraceItem{\TraceCode{ad4080\_\penalty150\relax{}csr\_\penalty150\relax{}power\_\penalty150\relax{}bridge} \href{https://github.com/defineiocc02/20bit\_SAR\_ADC\_Behaviour\_Verification/blob/feat/ad4080-peripherals/subsys/ad4080\_peripherals/interfaces/control\_review/ad4080\_csr\_power\_bridge.sv\#L24}{源码}}
\TraceField{手册规定}{\href{https://www.analog.com/media/en/technical-documentation/data-sheets/ad4080.pdf\#page=27}{p.27 Software Reset 章节}、\href{https://www.analog.com/media/en/technical-documentation/data-sheets/ad4080.pdf\#page=76}{p.76 表 34} 和 \href{https://www.analog.com/media/en/technical-documentation/data-sheets/ad4080.pdf\#page=83}{p.83 表 47} 给出模式、复位和 POR\_\penalty150\relax{}ANA\_\penalty150\relax{}FLAG 写 1 清零。}
\TraceField{代码、实测与尚缺的证据}{独立桥经过 0x02 请求、0x14 清位与新事件碰撞、非法 GPIO 拒绝。未改的 ADI no-OS C 驱动在独立装配跑过 47 项断言、126 帧 SPI。产品的 SPI 到原寄存器，电源请求另走顶层输入。\href{https://github.com/defineiocc02/20bit\_SAR\_ADC\_Behaviour\_Verification/blob/feat/ad4080-peripherals/subsys/ad4080\_peripherals/interfaces/control\_review/noos\_validation.txt}{no-OS 用例}}
\TraceField{裁定}{部分认可}
\end{TraceCard}
\begin{TraceCard}
\TraceItem{\TraceCode{ad4080\_\penalty150\relax{}top} \href{https://github.com/defineiocc02/20bit\_SAR\_ADC\_Behaviour\_Verification/blob/7408edf8db9a41eebdd756378b69d2775f107dec/subsys/ad4080\_dsp/rtl/ad4080\_top.sv\#L22}{源码}}
\TraceField{手册规定}{\href{https://www.analog.com/media/en/technical-documentation/data-sheets/ad4080.pdf\#page=33}{p.33 图 51、52} 划分配置 SPI 与结果端口，\href{https://www.analog.com/media/en/technical-documentation/data-sheets/ad4080.pdf\#page=54}{p.54 图 82} 和 \href{https://www.analog.com/media/en/technical-documentation/data-sheets/ad4080.pdf\#page=64}{p.64 图 96} 画出结果、事件和滤波的连接。}
\TraceField{代码、实测与尚缺的证据}{该模块串起 SPI、CSR、数字滤波、事件、GPIO 与捕获，输出仍是并行数字结果。产品装配临时引出外部结果存储的握手和地址；四银行存储由外围提供。\href{https://github.com/defineiocc02/20bit\_SAR\_ADC\_Behaviour\_Verification/blob/feat/ad4080-peripherals/subsys/ad4080\_peripherals/integration/product/bind\_dsp.py}{绑定器}}
\TraceField{裁定}{部分认可}
\end{TraceCard}
\begin{TraceCard}
\TraceItem{\TraceCode{ad4080\_\penalty150\relax{}peripheral\_\penalty150\relax{}product} \href{https://github.com/defineiocc02/20bit\_SAR\_ADC\_Behaviour\_Verification/blob/feat/ad4080-peripherals/subsys/ad4080\_peripherals/integration/product/ad4080\_peripheral\_product.sv\#L16}{源码}}
\TraceField{手册规定}{\href{https://www.analog.com/media/en/technical-documentation/data-sheets/ad4080.pdf\#page=33}{p.33 图 51、52}、\href{https://www.analog.com/media/en/technical-documentation/data-sheets/ad4080.pdf\#page=57}{p.57 表 27} 给出可配置接口以及捕获/读出模式。}
\TraceField{代码、实测与尚缺的证据}{已经打通补码输入\ensuremath{\rightarrow}配置 SPI\ensuremath{\rightarrow}原 DSP\ensuremath{\rightarrow}同一份 16K×20 存储\ensuremath{\rightarrow}单路 LVDS 逻辑位对，实测完整三字记录 46 对。当前记录组帧固定单路；未把实时输出、结果 SPI、AON 控制桥或受控起录器接成全产品。\href{https://github.com/defineiocc02/20bit\_SAR\_ADC\_Behaviour\_Verification/blob/feat/ad4080-peripherals/subsys/ad4080\_peripherals/integration/product/evidence/verification.txt}{真实产品结构和记录}}
\TraceField{裁定}{部分认可}
\end{TraceCard}


\hypertarget{trace-anchor-5}{}\pdfbookmark[0]{4. FIFO、结果存储与记录会话}{trace-heading-5}\TraceSection{4. FIFO、结果存储与记录会话}

\TraceTableHeading{模块或条目}
\begin{TraceCard}
\TraceItem{\TraceCode{fifo\_\penalty150\relax{}capture} \href{https://github.com/defineiocc02/20bit\_SAR\_ADC\_Behaviour\_Verification/blob/7408edf8db9a41eebdd756378b69d2775f107dec/subsys/ad4080\_dsp/rtl/fifo\_capture.sv\#L22}{源码}}
\TraceField{手册规定}{\href{https://www.analog.com/media/en/technical-documentation/data-sheets/ad4080.pdf\#page=57}{p.57 表 27} 列四种模式；\href{https://www.analog.com/media/en/technical-documentation/data-sheets/ad4080.pdf\#page=91}{p.91 表 56} 给水位，\href{https://www.analog.com/media/en/technical-documentation/data-sheets/ad4080.pdf\#page=63}{p.63 图 93–95} 举事件前后全记录例。}
\TraceField{代码、实测与尚缺的证据}{00 禁 FIFO，01 立即采集，10 事件后读 W 个，11 保留历史并读全深度。模式 11 有满史环写、事件地址 0x3FFB 和首两地址 0x3FFF\ensuremath{\rightarrow}0 的产品用例；随后主动中止，该用例没有完整读完 16,384 字。FULL/\penalty150\relax{}READ\_\penalty150\relax{}DONE 对真实移位寄存器的精确时刻仍需核对。\href{https://github.com/defineiocc02/20bit\_SAR\_ADC\_Behaviour\_Verification/blob/feat/ad4080-peripherals/subsys/ad4080\_peripherals/integration/product/evidence/verification.txt}{产品回执}}
\TraceField{裁定}{部分认可}
\end{TraceCard}
\begin{TraceCard}
\TraceItem{\TraceCode{result\_\penalty150\relax{}fifo} \href{https://github.com/defineiocc02/20bit\_SAR\_ADC\_Behaviour\_Verification/blob/7408edf8db9a41eebdd756378b69d2775f107dec/subsys/ad4080\_dsp/rtl/result\_fifo.sv\#L22}{源码}}
\TraceField{手册规定}{\href{https://www.analog.com/media/en/technical-documentation/data-sheets/ad4080.pdf\#page=56}{p.56 章节 Result FIFO} 明说单口、一次最多 16,384 个结果、采完再读；\href{https://www.analog.com/media/en/technical-documentation/data-sheets/ad4080.pdf\#page=57}{p.57 表 27} 定义使用方式。}
\TraceField{代码、实测与尚缺的证据}{原队列叶子允许普通读写同拍；在产品中移去内部数组，保留计数和指针，由捕获状态与物理握手互斥访问。裸队列不可直接当成器件单口证据。\href{https://github.com/defineiocc02/20bit\_SAR\_ADC\_Behaviour\_Verification/blob/feat/ad4080-peripherals/subsys/ad4080\_peripherals/verify/pinned\_fifo\_contract\_tb.sv}{全深度同拍检查} · \href{https://github.com/defineiocc02/20bit\_SAR\_ADC\_Behaviour\_Verification/blob/feat/ad4080-peripherals/subsys/ad4080\_peripherals/integration/product/bind\_dsp.py}{存储绑定器}}
\TraceField{裁定}{部分认可}
\end{TraceCard}
\begin{TraceCard}
\TraceItem{\TraceCode{result\_\penalty150\relax{}store} \href{https://github.com/defineiocc02/20bit\_SAR\_ADC\_Behaviour\_Verification/blob/feat/ad4080-peripherals/subsys/ad4080\_peripherals/rtl/result\_store.sv\#L12}{源码}}
\TraceField{手册规定}{\href{https://www.analog.com/media/en/technical-documentation/data-sheets/ad4080.pdf\#page=56}{p.56 Result FIFO 章节} 规定单端口 16K 结果缓冲，\href{https://www.analog.com/media/en/technical-documentation/data-sheets/ad4080.pdf\#page=20}{p.20 表 7} 给每笔 20 位。}
\TraceField{代码、实测与尚缺的证据}{物理端口每拍接收一笔读或一笔写，49,156 个检查周期遍历地址并核返回；产品另实写 16,384 笔。认可逻辑容量、单端口和数据保真。宏引脚、面积及布局需用目标工艺验证。\href{https://github.com/defineiocc02/20bit\_SAR\_ADC\_Behaviour\_Verification/blob/feat/ad4080-peripherals/subsys/ad4080\_peripherals/evidence/result-store.json}{存储回执} · \href{https://github.com/defineiocc02/20bit\_SAR\_ADC\_Behaviour\_Verification/blob/feat/ad4080-peripherals/subsys/ad4080\_peripherals/integration/product/evidence/verification.txt}{产品实测}}
\TraceField{裁定}{可认可}
\end{TraceCard}
\begin{TraceCard}
\TraceItem{\TraceCode{result\_\penalty150\relax{}bank\_\penalty150\relax{}4k20} \href{https://github.com/defineiocc02/20bit\_SAR\_ADC\_Behaviour\_Verification/blob/feat/ad4080-peripherals/subsys/ad4080\_peripherals/rtl/result\_store.sv\#L55}{源码}}
\TraceField{手册规定}{\href{https://www.analog.com/media/en/technical-documentation/data-sheets/ad4080.pdf\#page=56}{p.56 Result FIFO 章节} 公开总深度和单口访问。}
\TraceField{代码、实测与尚缺的证据}{四个 4096×20 银行、高两位片选、低 12 位字址、同步单拍返回是工程划分。Yosys 所见四银行可以验证本项目的组织，不能据此识别商品芯片中的银行数。\href{https://github.com/defineiocc02/20bit\_SAR\_ADC\_Behaviour\_Verification/blob/feat/ad4080-peripherals/subsys/ad4080\_peripherals/integration/product/evidence/verification.txt}{结构实测}}
\TraceField{裁定}{工程自定}
\end{TraceCard}
\begin{TraceCard}
\TraceItem{\TraceCode{ad4080\_\penalty150\relax{}product\_\penalty150\relax{}record\_\penalty150\relax{}session} \href{https://github.com/defineiocc02/20bit\_SAR\_ADC\_Behaviour\_Verification/blob/feat/ad4080-peripherals/subsys/ad4080\_peripherals/integration/product/ad4080\_product\_record\_session.sv\#L16}{源码}}
\TraceField{手册规定}{\href{https://www.analog.com/media/en/technical-documentation/data-sheets/ad4080.pdf\#page=59}{p.59 图 87、88} 和 \href{https://www.analog.com/media/en/technical-documentation/data-sheets/ad4080.pdf\#page=60}{p.60 图 89} 给出从满记录到读出的外部时序。}
\TraceField{代码、实测与尚缺的证据}{本模块在核心和 I/O 之间交接固定长度、控制锁与中止，末尾另发 record\_\penalty150\relax{}done\_\penalty150\relax{}core\_\penalty150\relax{}o。内部请求、掉电退场及双向同步是本工程协议。首/末 SRAM 读受理、字进入输出移位寄存器、最后一位送出是三个不同观察点。\href{https://github.com/defineiocc02/20bit\_SAR\_ADC\_Behaviour\_Verification/blob/feat/ad4080-peripherals/subsys/ad4080\_peripherals/integration/product/README.md}{记录会话实测}}
\TraceField{裁定}{工程自定}
\end{TraceCard}


\hypertarget{trace-anchor-6}{}\pdfbookmark[0]{5. 连续输出、采后组帧与 CNV}{trace-heading-6}\TraceSection{5. 连续输出、采后组帧与 CNV}

\TraceTableHeading{模块或条目}
\begin{TraceCard}
\TraceItem{\TraceCode{ad4080\_\penalty150\relax{}ddr\_\penalty150\relax{}edge\_\penalty150\relax{}words} \href{https://github.com/defineiocc02/20bit\_SAR\_ADC\_Behaviour\_Verification/blob/feat/ad4080-peripherals/subsys/ad4080\_peripherals/io/ad4080\_ddr\_edge\_words.sv\#L20}{源码}}
\TraceField{手册规定}{\href{https://www.analog.com/media/en/technical-documentation/data-sheets/ad4080.pdf\#page=46}{p.46 图 74、75}、\href{https://www.analog.com/media/en/technical-documentation/data-sheets/ad4080.pdf\#page=47}{p.47 图 76、77} 列出连续结果的单/双 lane 位序；\href{https://www.analog.com/media/en/technical-documentation/data-sheets/ad4080.pdf\#page=84}{p.84 表 48 的 INTF\_\penalty150\relax{}CHK\_\penalty150\relax{}EN} 给通信检查功能。}
\TraceField{代码、实测与尚缺的证据}{认可给定样本和 lane 配置下的数字位序。1,109 次向量检查、组合性质及原版 ADI FPGA 接收逻辑均核过；固定字 0xAC5D6。ODDR、DCO 相位、驱动能力和封装时序均在这一级之外。\href{https://github.com/defineiocc02/20bit\_SAR\_ADC\_Behaviour\_Verification/blob/feat/ad4080-peripherals/subsys/ad4080\_peripherals/io/handoff\_manifest.json}{发送回执} · \href{https://github.com/defineiocc02/20bit\_SAR\_ADC\_Behaviour\_Verification/blob/feat/ad4080-peripherals/subsys/ad4080\_peripherals/io/adi\_receiver/evidence.txt}{接收回执}}
\TraceField{裁定}{可认可}
\end{TraceCard}
\begin{TraceCard}
\TraceItem{\TraceCode{ad4080\_\penalty150\relax{}result\_\penalty150\relax{}transport} \href{https://github.com/defineiocc02/20bit\_SAR\_ADC\_Behaviour\_Verification/blob/feat/ad4080-peripherals/subsys/ad4080\_peripherals/transport/ad4080\_result\_transport.sv\#L23}{源码}}
\TraceField{手册规定}{\href{https://www.analog.com/media/en/technical-documentation/data-sheets/ad4080.pdf\#page=7}{p.7 表 2} 和 \href{https://www.analog.com/media/en/technical-documentation/data-sheets/ad4080.pdf\#page=46}{p.46 图 74、75} 给采样、输出时钟和连续读数的接口条件。}
\TraceField{代码、实测与尚缺的证据}{理想时钟下每档连续送 512 个码，单 lane 40/400 MHz、双 lane 40/200 MHz。Gray 指针、8×22 弹性队列和复位冲刷是我们选的实现。单/双 lane 空队列实测从并行码受理到首个位对分别为 7.503 ns 和 17.503 ns；这些数字的起点在 CNV 之后；产品记录口只选单 lane。\href{https://github.com/defineiocc02/20bit\_SAR\_ADC\_Behaviour\_Verification/blob/feat/ad4080-peripherals/subsys/ad4080\_peripherals/transport/verification.txt}{吞吐和首字实测}}
\TraceField{裁定}{部分认可}
\end{TraceCard}
\begin{TraceCard}
\TraceItem{\TraceCode{ad4080\_\penalty150\relax{}fifo\_\penalty150\relax{}record\_\penalty150\relax{}framer} \href{https://github.com/defineiocc02/20bit\_SAR\_ADC\_Behaviour\_Verification/blob/feat/ad4080-peripherals/subsys/ad4080\_peripherals/io/ad4080\_fifo\_record\_framer.sv\#L21}{源码}}
\TraceField{手册规定}{\href{https://www.analog.com/media/en/technical-documentation/data-sheets/ad4080.pdf\#page=59}{p.59 图 87、88} 画 SPI 单/四线、头和符号扩展；\href{https://www.analog.com/media/en/technical-documentation/data-sheets/ad4080.pdf\#page=60}{p.60 图 89} 画采后单路 LVDS 的头、20 位结果和 24 位尾。}
\TraceField{代码、实测与尚缺的证据}{SPI1 每记录一个 0xAA，SPI4 每条线各发 0xAA，载荷按 24 位符号扩展；台架核过图示前缀。单 LVDS 帧头和载荷已跑通，实际器件 CRC24 参数与覆盖区未闭合，要求外部 CRC 服务。台架只累计有效位对，等待下一字或等尾部时可能空拍，连续引脚时钟的预取及尾部截止时限另验。\href{https://github.com/defineiocc02/20bit\_SAR\_ADC\_Behaviour\_Verification/blob/feat/ad4080-peripherals/subsys/ad4080\_peripherals/io/evidence/record-source-check.md}{图示与读长裁定} · \href{https://github.com/defineiocc02/20bit\_SAR\_ADC\_Behaviour\_Verification/blob/feat/ad4080-peripherals/subsys/ad4080\_peripherals/io/FIFO\_RECORD.md}{组帧回执}}
\TraceField{裁定}{部分认可}
\end{TraceCard}
\begin{TraceCard}
\TraceItem{\TraceCode{ad4080\_\penalty150\relax{}result\_\penalty150\relax{}readback\_\penalty150\relax{}link} \href{https://github.com/defineiocc02/20bit\_SAR\_ADC\_Behaviour\_Verification/blob/feat/ad4080-peripherals/subsys/ad4080\_peripherals/integration/ad4080\_result\_readback\_link.sv\#L21}{源码}}
\TraceField{手册规定}{\href{https://www.analog.com/media/en/technical-documentation/data-sheets/ad4080.pdf\#page=56}{p.56 Result FIFO 章节}、\href{https://www.analog.com/media/en/technical-documentation/data-sheets/ad4080.pdf\#page=60}{p.60 图 89} 规定取出记录和单路 LVDS 载荷次序。}
\TraceField{代码、实测与尚缺的证据}{从唯一 SRAM 的真实有效读响应跨域，先后重组两个样本，核拒绝额外读。此叶子负责载荷、容量和响应交接；完整头/尾由外层负责。数据信用、8 字 FIFO、冲刷细节无厂商规定。\href{https://github.com/defineiocc02/20bit\_SAR\_ADC\_Behaviour\_Verification/blob/feat/ad4080-peripherals/subsys/ad4080\_peripherals/integration/evidence/readback\_link.txt}{线上重组证据}}
\TraceField{裁定}{部分认可}
\end{TraceCard}
\begin{TraceCard}
\TraceItem{\TraceCode{ad4080\_\penalty150\relax{}fifo\_\penalty150\relax{}lvds\_\penalty150\relax{}record\_\penalty150\relax{}link} \href{https://github.com/defineiocc02/20bit\_SAR\_ADC\_Behaviour\_Verification/blob/feat/ad4080-peripherals/subsys/ad4080\_peripherals/integration/ad4080\_fifo\_lvds\_record\_link.sv\#L20}{源码}}
\TraceField{手册规定}{\href{https://www.analog.com/media/en/technical-documentation/data-sheets/ad4080.pdf\#page=60}{p.60 图 89} 给出单路 LVDS 采后记录，\href{https://www.analog.com/media/en/technical-documentation/data-sheets/ad4080.pdf\#page=51}{p.51 表 22、23} 单列结果 SPI 的信号与数据线序。}
\TraceField{代码、实测与尚缺的证据}{该联合入口将组帧器固定为单 LVDS，SPI 样本口接零。独立接收器得到两样本、36 对的工程 CRC 记录；产品得到三样本、46 对。SPI1/4 组帧台架和实时双 lane 台架不能证明这些端口进入当前产品。\href{https://github.com/defineiocc02/20bit\_SAR\_ADC\_Behaviour\_Verification/blob/feat/ad4080-peripherals/subsys/ad4080\_peripherals/integration/evidence/readback\_link.txt}{联合回执} · \href{https://github.com/defineiocc02/20bit\_SAR\_ADC\_Behaviour\_Verification/blob/feat/ad4080-peripherals/subsys/ad4080\_peripherals/integration/product/evidence/verification.txt}{产品回执}}
\TraceField{裁定}{部分认可}
\end{TraceCard}
\begin{TraceCard}
\TraceItem{\TraceCode{cnv\_\penalty150\relax{}ctrl} \href{https://github.com/defineiocc02/20bit\_SAR\_ADC\_Behaviour\_Verification/blob/7408edf8db9a41eebdd756378b69d2775f107dec/subsys/ad4080\_dsp/rtl/cnv\_ctrl.sv\#L22}{源码}}
\TraceField{手册规定}{\href{https://www.analog.com/media/en/technical-documentation/data-sheets/ad4080.pdf\#page=7}{p.7 表 2} 给出 tCONV、CNV/CLK 对齐；\href{https://www.analog.com/media/en/technical-documentation/data-sheets/ad4080.pdf\#page=83}{p.83 表 47 的 ADC\_\penalty150\relax{}CNV\_\penalty150\relax{}ERR} 表示违反最小转换时序时数据无效。}
\TraceField{代码、实测与尚缺的证据}{原独立模块仅同步已译码 CNV，检测逻辑上升沿碰到 busy。没有测量 tCONV 的计时器；生产外围未实例化它，也没有将故障送入 0x14[2]。逻辑 busy 冲突无法证明实现了器件的时序故障检测。\href{https://github.com/defineiocc02/20bit\_SAR\_ADC\_Behaviour\_Verification/blob/feat/ad4080-peripherals/subsys/ad4080\_peripherals/README.md}{DSP 入口说明}}
\TraceField{裁定}{不认可}
\end{TraceCard}


\hypertarget{trace-anchor-7}{}\pdfbookmark[0]{6. 电源、上电装载与自测}{trace-heading-7}\TraceSection{6. 电源、上电装载与自测}

\TraceTableHeading{模块或条目}
\begin{TraceCard}
\TraceItem{\TraceCode{ad4080\_\penalty150\relax{}pwr\_\penalty150\relax{}supervisor} \href{https://github.com/defineiocc02/20bit\_SAR\_ADC\_Behaviour\_Verification/blob/feat/ad4080-peripherals/subsys/ad4080\_peripherals/power/ad4080\_pwr\_supervisor.sv\#L21}{源码}}
\TraceField{手册规定}{\href{https://www.analog.com/media/en/technical-documentation/data-sheets/ad4080.pdf\#page=24}{p.24 图 40} 区分核心 1.1 V 与 REFIN 两个监测；\href{https://www.analog.com/media/en/technical-documentation/data-sheets/ad4080.pdf\#page=26}{p.26 图 43、表 11}、\href{https://www.analog.com/media/en/technical-documentation/data-sheets/ad4080.pdf\#page=27}{p.27 图 44} 给模式及恢复，\href{https://www.analog.com/media/en/technical-documentation/data-sheets/ad4080.pdf\#page=83}{p.83 表 47} 给 W1C 状态。}
\TraceField{代码、实测与尚缺的证据}{在逻辑 PG/故障输入下测分域许可、停钟、模式和复位，AON 计时器按配置频率向上取整。手册的 300 µs 初始化、待机/睡眠唤醒及额外精度恢复各有起点。未测模拟阈值、公差、唤醒精度或真实 LDO；独立控制用例与产品的 SPI 命令路径另计。\href{https://github.com/defineiocc02/20bit\_SAR\_ADC\_Behaviour\_Verification/blob/feat/ad4080-peripherals/subsys/ad4080\_peripherals/power/verification.txt}{逻辑电源回执}}
\TraceField{裁定}{部分认可}
\end{TraceCard}
\begin{TraceCard}
\TraceItem{\TraceCode{ad4080\_\penalty150\relax{}pwr\_\penalty150\relax{}reset\_\penalty150\relax{}sync} \href{https://github.com/defineiocc02/20bit\_SAR\_ADC\_Behaviour\_Verification/blob/feat/ad4080-peripherals/subsys/ad4080\_peripherals/power/ad4080\_pwr\_reset\_sync.sv\#L17}{源码}}
\TraceField{手册规定}{\href{https://www.analog.com/media/en/technical-documentation/data-sheets/ad4080.pdf\#page=24}{p.24 图 40}、\href{https://www.analog.com/media/en/technical-documentation/data-sheets/ad4080.pdf\#page=25}{p.25 表 10、图 41–42} 给域失电、上电次序和复位要求。}
\TraceField{代码、实测与尚缺的证据}{本工程采用异步封锁、在本时钟域同步释放的复位同步器；手册没有公开触发器级数。数字停钟测试不能代替目标库 MTBF、CDC/RDC、复位树和时序签核。\href{https://github.com/defineiocc02/20bit\_SAR\_ADC\_Behaviour\_Verification/blob/feat/ad4080-peripherals/subsys/ad4080\_peripherals/power/README.md}{电源和跨域验收范围}}
\TraceField{裁定}{工程自定}
\end{TraceCard}
\begin{TraceCard}
\TraceItem{\TraceCode{ad4080\_\penalty150\relax{}peripheral\_\penalty150\relax{}services} \href{https://github.com/defineiocc02/20bit\_SAR\_ADC\_Behaviour\_Verification/blob/feat/ad4080-peripherals/subsys/ad4080\_peripherals/integration/ad4080\_peripheral\_services.sv\#L22}{源码}}
\TraceField{手册规定}{\href{https://www.analog.com/media/en/technical-documentation/data-sheets/ad4080.pdf\#page=24}{p.24 图 40}、\href{https://www.analog.com/media/en/technical-documentation/data-sheets/ad4080.pdf\#page=26}{p.26 表 11}、\href{https://www.analog.com/media/en/technical-documentation/data-sheets/ad4080.pdf\#page=56}{p.56 Result FIFO 章节} 给电源资格与缓冲访问的外部约束。}
\TraceField{代码、实测与尚缺的证据}{服务壳协调供电、初始化、记录、存储自测和退场。验证了未完采集阻挡、授权前不破坏及掉电恢复。商业芯片的内部仲裁没有公开。产品中 SPI 模式与软复位请求仍同 AON 管理口分开，不能宣称全片低功耗流程通过。\href{https://github.com/defineiocc02/20bit\_SAR\_ADC\_Behaviour\_Verification/blob/feat/ad4080-peripherals/subsys/ad4080\_peripherals/integration/evidence/regression.txt}{13 个服务场景}}
\TraceField{裁定}{部分认可}
\end{TraceCard}
\begin{TraceCard}
\TraceItem{\TraceCode{otp\_\penalty150\relax{}shadow\_\penalty150\relax{}loader} \href{https://github.com/defineiocc02/20bit\_SAR\_ADC\_Behaviour\_Verification/blob/feat/ad4080-peripherals/subsys/ad4080\_peripherals/otp/otp\_shadow\_loader.sv\#L20}{源码}}
\TraceField{手册规定}{\href{https://www.analog.com/media/en/technical-documentation/data-sheets/ad4080.pdf\#page=39}{p.39 表 16} 与 \href{https://www.analog.com/media/en/technical-documentation/data-sheets/ad4080.pdf\#page=83}{p.83 表 47} 只给 ROM CRC/ECC 错误读位；\href{https://www.analog.com/media/en/technical-documentation/data-sheets/ad4080.pdf\#page=71}{p.71 系统校正章节} 确认出厂校准。}
\TraceField{代码、实测与尚缺的证据}{工程装载器经外部宏读取，待外部校验器同意才提交影子值；未请先答等场景有定向测试。手册没有确定一次性可编程宏、镜像尺寸、地址、编程电压、校验算法或者影子提交。当前产品还没有把校验失败送至 0x14[1]。\href{https://github.com/defineiocc02/20bit\_SAR\_ADC\_Behaviour\_Verification/blob/feat/ad4080-peripherals/subsys/ad4080\_peripherals/otp/verification.txt}{装载器回执}}
\TraceField{裁定}{工程自定}
\end{TraceCard}
\begin{TraceCard}
\TraceItem{\TraceCode{result\_\penalty150\relax{}store\_\penalty150\relax{}mbist} \href{https://github.com/defineiocc02/20bit\_SAR\_ADC\_Behaviour\_Verification/blob/feat/ad4080-peripherals/subsys/ad4080\_peripherals/memory/result\_store\_mbist.sv\#L18}{源码}}
\TraceField{手册规定}{\href{https://www.analog.com/media/en/technical-documentation/data-sheets/ad4080.pdf\#page=56}{p.56 Result FIFO 章节} 给单口结果存储；\href{https://www.analog.com/media/en/technical-documentation/data-sheets/ad4080.pdf\#page=83}{p.83 表 47} 的故障位说的是 ROM CRC/ECC。}
\TraceField{代码、实测与尚缺的证据}{March C−、测试底图、独占和重建策略由工程选定。两种底图各执行 163,840 次访问，检出了注入故障。手册没有披露 SRAM 测试算法、修复、扫描链或 ATE 覆盖率。\href{https://github.com/defineiocc02/20bit\_SAR\_ADC\_Behaviour\_Verification/blob/feat/ad4080-peripherals/subsys/ad4080\_peripherals/memory/evidence.json}{MBIST 实测}}
\TraceField{裁定}{工程自定}
\end{TraceCard}


\hypertarget{trace-anchor-8}{}\pdfbookmark[0]{7. 寄存器和时序中容易误认的部分}{trace-heading-8}\TraceSection{7. 寄存器和时序中容易误认的部分}

\TraceTableHeading{对外事实}
\begin{TraceCard}
\TraceItem{系统校正与传输延迟}
\TraceField{手册精确位置}{\href{https://www.analog.com/media/en/technical-documentation/data-sheets/ad4080.pdf\#page=20}{p.20 图 33、表 7}、\href{https://www.analog.com/media/en/technical-documentation/data-sheets/ad4080.pdf\#page=48}{p.48 ADC Result Latency 章节}、\href{https://www.analog.com/media/en/technical-documentation/data-sheets/ad4080.pdf\#page=71}{p.71 式 (2)}、\href{https://www.analog.com/media/en/technical-documentation/data-sheets/ad4080.pdf\#page=93}{p.93 表 61}}
\TraceField{本次采用的解释及产品限制}{VREF=3 V 时差分 1 LSB 为 6 V/2\textsuperscript{2}\textsuperscript{0}；OFFSET 是 12 位补码，GAIN 是 10 位直码。0x001 和 0x3FF 对应约 −/+1.5594\%，0x200 为单位增益。0x000 与各档精确位级舍入未列出。40 MS/s 单 lane 示例有 46.25 ns（单位增益）和 51.25 ns（启用增益校正）两种读数，不能用模型的两拍延迟代替。}
\end{TraceCard}
\begin{TraceCard}
\TraceItem{OFFSET 起址、字节序}
\TraceField{手册精确位置}{\href{https://www.analog.com/media/en/technical-documentation/data-sheets/ad4080.pdf\#page=71}{p.71 Offset Error Correction 章节}、\href{https://www.analog.com/media/en/technical-documentation/data-sheets/ad4080.pdf\#page=74}{p.74 表 31}、\href{https://www.analog.com/media/en/technical-documentation/data-sheets/ad4080.pdf\#page=92}{p.92 表 59、60}；\href{https://www.analog.com/media/en/technical-documentation/data-sheets/ad4080.pdf\#page=36}{p.36 表 15}、\href{https://www.analog.com/media/en/technical-documentation/data-sheets/ad4080.pdf\#page=38}{p.38 图 57–59}}
\TraceField{本次采用的解释及产品限制}{p.71 文字写 0x24/25，但 0x24 已给低阈值高字节。地址表、图 138 和表 60 支持 0x25/26，GAIN 为 0x27/28。升序从低地址/低字节起，降序从高地址/高字节起。p.38 与 p.41 个别文字对起址有反向描述，按图核验整实体并读回。}
\end{TraceCard}
\begin{TraceCard}
\TraceItem{身份、普通访问、配置 CRC}
\TraceField{手册精确位置}{\href{https://www.analog.com/media/en/technical-documentation/data-sheets/ad4080.pdf\#page=35}{p.35 Write Access 章节}、\href{https://www.analog.com/media/en/technical-documentation/data-sheets/ad4080.pdf\#page=41}{p.41 表 18、19}、\href{https://www.analog.com/media/en/technical-documentation/data-sheets/ad4080.pdf\#page=77}{p.77 表 35、36}、\href{https://www.analog.com/media/en/technical-documentation/data-sheets/ad4080.pdf\#page=79}{p.79 表 41、42}、\href{https://www.analog.com/media/en/technical-documentation/data-sheets/ad4080.pdf\#page=80}{p.80 表 43}}
\TraceField{本次采用的解释及产品限制}{独立 no-OS 夹具识别产品 0x0050，p.77 的 0x04 页首写复位 0x00、同页表 36 却写 0x50；当前产品原参数未绑定。VENDOR 低/高字节为 0x56/\penalty150\relax{}0x04。配置 CRC-8 的多项式为 x\textsuperscript{8}+x\textsuperscript{2}+x+1，有明确资料；FIFO CRC24 为另一问题。基本升降址流式事务与 0x0E LOOP\_\penalty150\relax{}COUNT 是两种功能，该页还将后者标为不支持。}
\end{TraceCard}
\begin{TraceCard}
\TraceItem{模式、状态与输出控制}
\TraceField{手册精确位置}{\href{https://www.analog.com/media/en/technical-documentation/data-sheets/ad4080.pdf\#page=76}{p.76 表 34}、\href{https://www.analog.com/media/en/technical-documentation/data-sheets/ad4080.pdf\#page=83}{p.83 表 47}、\href{https://www.analog.com/media/en/technical-documentation/data-sheets/ad4080.pdf\#page=84}{p.84 表 48}、\href{https://www.analog.com/media/en/technical-documentation/data-sheets/ad4080.pdf\#page=85}{p.85 表 49、50}、\href{https://www.analog.com/media/en/technical-documentation/data-sheets/ad4080.pdf\#page=86}{p.86 表 51}}
\TraceField{本次采用的解释及产品限制}{0x02 操作模式与 AON 请求缺产品级贯通。0x14[3] 模拟 POR、[2] CNV 错误和 [1] ROM 错误缺同源故障输入。0x15–17 的模式、lane、相位/自时钟等不能从实际产品 CSR 配置；0x18 电源控制字段亦不可凭电源主管台架认可。}
\end{TraceCard}
\begin{TraceCard}
\TraceItem{事件、水位与状态}
\TraceField{手册精确位置}{\href{https://www.analog.com/media/en/technical-documentation/data-sheets/ad4080.pdf\#page=53}{p.53 表 25}、\href{https://www.analog.com/media/en/technical-documentation/data-sheets/ad4080.pdf\#page=55}{p.55 图 83}、\href{https://www.analog.com/media/en/technical-documentation/data-sheets/ad4080.pdf\#page=57}{p.57 表 27}、\href{https://www.analog.com/media/en/technical-documentation/data-sheets/ad4080.pdf\#page=59}{p.59 图 87、88}、\href{https://www.analog.com/media/en/technical-documentation/data-sheets/ad4080.pdf\#page=63}{p.63 图 93–95}、\href{https://www.analog.com/media/en/technical-documentation/data-sheets/ad4080.pdf\#page=91}{p.91 表 56}}
\TraceField{本次采用的解释及产品限制}{0x1D/1E 默认 0x4000。模式 11 的 W 为 4 的倍数，捕获事件位置会回写水位实体；图 95 的计数器没有独立 EVENT\_\penalty150\relax{}COUNT 地址。在 W=8192 例子里图 93 和图 94 对首个事件后样本标了 8193 与 8192，片内物理地址和读回样本序号应分别核对。表 25 的 READ\_\penalty150\relax{}DONE 描述与时序图存在置/清歧义，产品区分 SRAM 读受理和现场线位完成。}
\end{TraceCard}
\begin{TraceCard}
\TraceItem{滤波建立和首字}
\TraceField{手册精确位置}{\href{https://www.analog.com/media/en/technical-documentation/data-sheets/ad4080.pdf\#page=66}{p.66 图 97、表 29}、\href{https://www.analog.com/media/en/technical-documentation/data-sheets/ad4080.pdf\#page=67}{p.67 图 98}、\href{https://www.analog.com/media/en/technical-documentation/data-sheets/ad4080.pdf\#page=68}{p.68 图 99}、\href{https://www.analog.com/media/en/technical-documentation/data-sheets/ad4080.pdf\#page=69}{p.69 Digital Filter Conversion Pulses 章节}、\href{https://www.analog.com/media/en/technical-documentation/data-sheets/ad4080.pdf\#page=94}{p.94 表 62}}
\TraceField{本次采用的解释及产品限制}{N 为 sinc 抽取倍率。sinc1 的建立计数为 2+N+1，sinc5 为 2+5N+4。补偿路径有额外 2 倍抽取，表 29 按选项列出实际总倍率。图 99 区分低有效 RDY、写 RAM 与 FULL；N=2、W=3 图中从 E1 同步至 E10 FULL，p.69 的 W·D+建立式按字面得 11，观测起点须由器件波形确认。}
\end{TraceCard}

\textbf{默认配置和实验配置不同。} 产品校准用例显式绑定 OFFSET 的最低地址 0x25、降序多字节、K=1、SHIFT=15；常规 SINC1 产品例使用 SHIFT=14。产品默认不开放未绑定的非单位增益和补偿系数。普通配置写 SDO 已按原文改为驱动 0，但旧有数值链与图 96 的次序差异仍在。独立 no-OS 夹具的产品 ID 与原驱动可读写性也不能移记到这个产品入口。\href{https://github.com/defineiocc02/20bit\_SAR\_ADC\_Behaviour\_Verification/blob/feat/ad4080-peripherals/subsys/ad4080\_peripherals/integration/product/ad4080\_peripheral\_product.sv}{产品端口与绑定} · \href{https://github.com/defineiocc02/20bit\_SAR\_ADC\_Behaviour\_Verification/blob/feat/ad4080-peripherals/subsys/ad4080\_peripherals/integration/product/README.md}{校准与常规实测}\par


\hypertarget{trace-anchor-9}{}\pdfbookmark[0]{8. 哪些功能没有芯片手册依据}{trace-heading-9}\TraceSection{8. 哪些功能没有芯片手册依据}

\begin{itemize}
\item \textbf{结果 CRC。} 图 87、88 支持 SPI1/\penalty150\relax{}SPI4 图示访问前缀；图 89 画出单路 LVDS 的 24 位尾。没有获得与不同数据/长度对应的真实 CRC24 金样，也没有从 Rev. B 找到多项式、初值、反射、终异或及完整覆盖范围。工程余数 0x71F3EC 只属于已声明的测试参数。ADI 主机固件和 FPGA 原源码可复核其读帧动作；离线 CRC 程序、合成例和理想原语属于测试工具。\href{https://www.analog.com/media/en/technical-documentation/data-sheets/ad4080.pdf\#page=59}{手册图 87、88} · \href{https://www.analog.com/media/en/technical-documentation/data-sheets/ad4080.pdf\#page=60}{手册图 89} · \href{https://github.com/defineiocc02/20bit\_SAR\_ADC\_Behaviour\_Verification/blob/feat/ad4080-peripherals/subsys/ad4080\_peripherals/io/adi\_crc/README.md}{CRC 采集与判别} · \href{https://github.com/defineiocc02/20bit\_SAR\_ADC\_Behaviour\_Verification/blob/feat/ad4080-peripherals/subsys/ad4080\_peripherals/io/adi\_receiver/README.md}{ADI FPGA 实测}\par
\item \textbf{电源和工厂存储。} 手册给出双 LDO 内供和外供接法、核心/参考电压监测、上电后 300 µs 初始化和出厂校准。待机 100 µs、睡眠 180 µs 的唤醒时间之后，图 44 另画没有定值的恢复时间；精度恢复不能用逻辑定时器盖章。查 95 页可检索正文和相关图，能定位只读 ROM CRC/ECC 位，找不到 OTP 工艺、启动镜像、March 算法、修复或扫描资源的器件披露。数据手册的 ROM、配置 CRC8、记录 CRC24 是三个对象。\href{https://www.analog.com/media/en/technical-documentation/data-sheets/ad4080.pdf\#page=23}{p.23 内供图 38} · \href{https://www.analog.com/media/en/technical-documentation/data-sheets/ad4080.pdf\#page=24}{p.24 外供和 POR 图 39、40} · \href{https://www.analog.com/media/en/technical-documentation/data-sheets/ad4080.pdf\#page=27}{图 44} · \href{https://www.analog.com/media/en/technical-documentation/data-sheets/ad4080.pdf\#page=83}{表 47}\par
\item \textbf{版图和面积。} Rev. B 给出 49 球 CSP BGA、引脚及外部供电连接；手册没有给 die 尺寸、工艺、各数字功能面积、小 SRAM 的用途、OTP 区大小或宏 LEF。规划 GDS、四银行排布、面积预算、PDN 通道与 SRAM 接收脚本都是项目交付。不能凭封装长宽反推出裸片占比。论文/讲稿原型的 Digital 框可以单列为原型量测，不能填为 AD4080 商品芯片模块面积。\href{https://www.analog.com/media/en/technical-documentation/data-sheets/ad4080.pdf\#page=10}{表 6 引脚} · \href{https://www.analog.com/media/en/technical-documentation/data-sheets/ad4080.pdf\#page=95}{封装页} · \href{https://github.com/defineiocc02/20bit\_SAR\_ADC\_Behaviour\_Verification/blob/feat/ad4080-peripherals/subsys/ad4080\_peripherals/budget/README.md}{项目面积和宏资料门}\par
\item \textbf{转换核与模拟宏。} 36 项职责中，片组调度、Flash/\penalty150\relax{}SAR 事件跟随、双 SADC、DEM 选元、RDAC、dither 生成/预置与扣除、单元权重训练与应用、余差路径、ADC2 量化及其归属共 12 项属于转换、模拟或相应的片外权重估计。Rev. B 图 31、32 支持 SAR 与连续采集框架，未披露本仓的选片、扰动、权重或内部端口细节。系统 OFFSET/\penalty150\relax{}GAIN 是另一个环节。登记在此处的转换核沿用既有功能组，审核结论止于交接边界。\href{https://www.analog.com/media/en/technical-documentation/data-sheets/ad4080.pdf\#page=19}{架构图 31、32} · \href{https://github.com/defineiocc02/20bit\_SAR\_ADC\_Behaviour\_Verification/blob/feat/ad4080-peripherals/subsys/ad4080\_peripherals/interfaces/function\_ownership.csv}{36 项职责逐条归属}\par
\item \textbf{签核和测量。} \TraceCode{4×4096×20}、CDC 深度、常开时钟上界、Pad 位置与面积是设计参数。要认可器件等效性，还需要实物在同一时基下的 CNV/\penalty150\relax{}CLK/\penalty150\relax{}GPIO/\penalty150\relax{}串行位流、增益与 OFFSET 联合扫描，以及供应方宏和 PVT 资料。布图、GF(2) 分析、no-\penalty150\relax{}OS/\penalty150\relax{}HDL 联测和既有 29 组回归各自回答不同的问题。\href{https://github.com/defineiocc02/20bit\_SAR\_ADC\_Behaviour\_Verification/blob/feat/ad4080-peripherals/subsys/ad4080\_peripherals/evidence/peripheral\_open\_items.csv}{物理与逻辑未闭合项} · \href{https://github.com/defineiocc02/20bit\_SAR\_ADC\_Behaviour\_Verification/blob/feat/ad4080-peripherals/subsys/ad4080\_peripherals/evidence/peripheral-review.json}{本项目回归回执}\par
\end{itemize}


\hypertarget{trace-anchor-10}{}\pdfbookmark[0]{9. 判据、待办与复核方法}{trace-heading-10}\TraceSection{9. 判据、待办与复核方法}

\textbf{需要分别处理的现场问题。}\par

\begin{enumerate}
\item \textbf{[需仓主决策] 校正单元的寄存器语义。} 片外实测增益表已按图 96 给出乘前 OFFSET 初值，工程 K/SHIFT 仍准确重现现用 RTL。若要求产品达到手册，应由该 DSP 的维护者在同一实现上统一 RTL 与整数参考，让 +1024 的判别向量过关，再用正、负 OFFSET 和多档 GAIN 做联合验证。板上绝对精度另需实测 VREF 及已标定输入。\href{https://www.analog.com/media/en/technical-documentation/data-sheets/ad4080.pdf\#page=64}{手册图 96} · \href{https://github.com/defineiocc02/20bit\_SAR\_ADC\_Behaviour\_Verification/blob/feat/ad4080-peripherals/subsys/ad4080\_peripherals/evidence/peripheral\_open\_items.csv}{开项与解除条件}\par
\item \textbf{[需仓主决策] 首先交付哪几种产品工作模式。} 若承诺全手册接口，需要把 0x15–17 的实际选择、连续单/双 lane、采后 SPI、滤波同步与有效就绪、配置到真实电源命令都接入同一个产品顶层，再用主机写寄存器、观察 GPIO、完整读到最后一位。外供 IOVDD 的待机读 FIFO 有专门条件；单路 LVDS 记录回归无法替代它。还应覆盖 CNV 停走、字间连续性与模式 11 的整记录读回。\href{https://www.analog.com/media/en/technical-documentation/data-sheets/ad4080.pdf\#page=26}{表 11} · \href{https://www.analog.com/media/en/technical-documentation/data-sheets/ad4080.pdf\#page=33}{数据接口图 51、52}\par
\item \textbf{核实事件语义与故障报告。} p.55 图 83 区分越阈后的活动电平 ALERT 和写 1 才清除的粘滞状态；p.56 图 84、85 对迟滞电平和路由信号另作描述。当前回窗后实得 HI\_DTCT=0、ALERT\_N=1、HI\_\penalty150\relax{}STATUS=1。RTL 对阈值相等采用高端 \ensuremath{\geq}、低端 \ensuremath{\leq}，W1C 同沿优先清除粘滞位，事件脉冲可仍然产生。须在同一 CNV 波形上分别采 GPIO、CSR 和 FIFO 触发。冷启动从未提供 REFIN 时测试字应停；先提供再移除时测试字可用。服务壳目前没有跟踪这种历史。\href{https://www.analog.com/media/en/technical-documentation/data-sheets/ad4080.pdf\#page=55}{图 83} · \href{https://www.analog.com/media/en/technical-documentation/data-sheets/ad4080.pdf\#page=56}{图 84、85} · \href{https://www.analog.com/media/en/technical-documentation/data-sheets/ad4080.pdf\#page=24}{冷启动参考条件} · \href{https://github.com/defineiocc02/20bit\_SAR\_ADC\_Behaviour\_Verification/blob/feat/ad4080-peripherals/subsys/ad4080\_peripherals/interfaces/control\_review/validation.txt}{回窗反例}\par
\item \textbf{[需仓主决策] 未披露的内部实现。} OTP 影子表、March C−、四银行 SRAM、跨域 FIFO、真实宏、floorplan 与测试覆盖可按本项目独立指标验收。若要标为 AD4080 器件实现，需要另外取得适用的直接证据。现有公开资料只支撑表中列出的外部行为。\href{https://www.analog.com/media/en/technical-documentation/data-sheets/ad4080.pdf\#page=1}{功能/接口图 1} · \href{https://github.com/defineiocc02/20bit\_SAR\_ADC\_Behaviour\_Verification/blob/feat/ad4080-peripherals/subsys/ad4080\_peripherals/interfaces/function\_ownership.csv}{工程功能矩阵}\par
\end{enumerate}

验证专用清单：\TraceCode{BUFIO}、\TraceCode{BUFR}、\TraceCode{IBUFDS}、\TraceCode{IBUFGDS}、\TraceCode{IDELAYCTRL}、\TraceCode{IDELAYE2}、\TraceCode{ISERDESE2}、\TraceCode{ad4080\_\penalty150\relax{}ddr\_\penalty150\relax{}edge\_\penalty150\relax{}words\_\penalty150\relax{}properties}、\TraceCode{ad4080\_\penalty150\relax{}ddr\_\penalty150\relax{}edge\_\penalty150\relax{}words\_\penalty150\relax{}tb}、\TraceCode{ad4080\_\penalty150\relax{}fifo\_\penalty150\relax{}record\_\penalty150\relax{}framer\_\penalty150\relax{}tb}、\TraceCode{adi\_\penalty150\relax{}receiver\_\penalty150\relax{}tb}、\TraceCode{calibration\_\penalty150\relax{}verify\_\penalty150\relax{}tb}、\TraceCode{control\_\penalty150\relax{}review\_\penalty150\relax{}tb}、\TraceCode{noos\_\penalty150\relax{}spi\_\penalty150\relax{}rtl}、\TraceCode{peripheral\_\penalty150\relax{}readback\_\penalty150\relax{}link\_\penalty150\relax{}tb}、\TraceCode{peripheral\_\penalty150\relax{}services\_\penalty150\relax{}tb}、\TraceCode{pinned\_\penalty150\relax{}fifo\_\penalty150\relax{}contract\_\penalty150\relax{}tb}、\TraceCode{power\_\penalty150\relax{}bridge\_\penalty150\relax{}tb}、\TraceCode{result\_\penalty150\relax{}store\_\penalty150\relax{}mbist\_\penalty150\relax{}tb}、\TraceCode{store\_\penalty150\relax{}tb}、\TraceCode{tb\_\penalty150\relax{}ad4080\_\penalty150\relax{}peripheral\_\penalty150\relax{}product}、\TraceCode{tb\_\penalty150\relax{}ad4080\_\penalty150\relax{}pwr\_\penalty150\relax{}supervisor}、\TraceCode{tb\_\penalty150\relax{}ad4080\_\penalty150\relax{}result\_\penalty150\relax{}transport}、\TraceCode{tb\_\penalty150\relax{}dsp\_\penalty150\relax{}path}、\TraceCode{tb\_\penalty150\relax{}otp\_\penalty150\relax{}shadow\_\penalty150\relax{}loader}。\par

从仓库根运行。第一条按声明和文件清单检查覆盖、源码目标及手册定位；第二条以器件一致性为条件，遇部分认可或不认可时预期退出 2。这个检查器能发现漏列模块和缺少出处，逐图阅读及行为裁定仍由审查者负责。\par

\begin{TraceCodeBlock}
\TraceCodeLine{p=\penalty150\relax{}subsys/\penalty150\relax{}ad4080\_\penalty150\relax{}peripherals}
\TraceCodeLine{node\hspace*{\fontdimen2\font}\allowbreak{}"\$p/\penalty150\relax{}verify/\penalty150\relax{}check\_\penalty150\relax{}manual\_\penalty150\relax{}trace.\penalty150\relax{}mjs"\hspace*{\fontdimen2\font}\allowbreak{}-\penalty150\relax{}-\penalty150\relax{}check-\penalty150\relax{}local}
\TraceCodeLine{node\hspace*{\fontdimen2\font}\allowbreak{}"\$p/\penalty150\relax{}verify/\penalty150\relax{}check\_\penalty150\relax{}manual\_\penalty150\relax{}trace.\penalty150\relax{}mjs"\hspace*{\fontdimen2\font}\allowbreak{}-\penalty150\relax{}-\penalty150\relax{}require-\penalty150\relax{}device-\penalty150\relax{}match}
\end{TraceCodeBlock}

做模块复测时，校准对应 \TraceCode{adctoolbox\_\penalty150\relax{}host\hspace*{\fontdimen2\font}\allowbreak{}/\penalty150\relax{}\hspace*{\fontdimen2\font}\allowbreak{}adctoolbox\_\penalty150\relax{}rtl\hspace*{\fontdimen2\font}\allowbreak{}/\penalty150\relax{}\hspace*{\fontdimen2\font}\allowbreak{}product\_\penalty150\relax{}calibration}；滤波、事件与 GPIO 对应 \TraceCode{dsp\_\penalty150\relax{}data\hspace*{\fontdimen2\font}\allowbreak{}/\penalty150\relax{}\hspace*{\fontdimen2\font}\allowbreak{}dsp\_\penalty150\relax{}control\hspace*{\fontdimen2\font}\allowbreak{}/\penalty150\relax{}\hspace*{\fontdimen2\font}\allowbreak{}adi\_\penalty150\relax{}noos}；FIFO 和实体存储对应 \TraceCode{store\hspace*{\fontdimen2\font}\allowbreak{}/\penalty150\relax{}\hspace*{\fontdimen2\font}\allowbreak{}queue\_\penalty150\relax{}port\_\penalty150\relax{}full\_\penalty150\relax{}depth\hspace*{\fontdimen2\font}\allowbreak{}/\penalty150\relax{}\hspace*{\fontdimen2\font}\allowbreak{}product\_\penalty150\relax{}binding}；输出接口对应 \TraceCode{edge\_\penalty150\relax{}words\hspace*{\fontdimen2\font}\allowbreak{}/\penalty150\relax{}\hspace*{\fontdimen2\font}\allowbreak{}record\_\penalty150\relax{}frame\_\penalty150\relax{}0\hspace*{\fontdimen2\font}\allowbreak{}/\penalty150\relax{}\hspace*{\fontdimen2\font}\allowbreak{}record\_\penalty150\relax{}link\hspace*{\fontdimen2\font}\allowbreak{}/\penalty150\relax{}\hspace*{\fontdimen2\font}\allowbreak{}transport\hspace*{\fontdimen2\font}\allowbreak{}/\penalty150\relax{}\hspace*{\fontdimen2\font}\allowbreak{}adi\_\penalty150\relax{}receiver}；启动与自测对应 \TraceCode{power\hspace*{\fontdimen2\font}\allowbreak{}/\penalty150\relax{}\hspace*{\fontdimen2\font}\allowbreak{}otp\hspace*{\fontdimen2\font}\allowbreak{}/\penalty150\relax{}\hspace*{\fontdimen2\font}\allowbreak{}mbist\hspace*{\fontdimen2\font}\allowbreak{}/\penalty150\relax{}\hspace*{\fontdimen2\font}\allowbreak{}services}。同一命令可以重复传 \TraceCode{-\penalty150\relax{}-\penalty150\relax{}case}。含官方互操作测试时，以 \TraceCode{-\penalty150\relax{}-\penalty150\relax{}fetch-\penalty150\relax{}adi} 获取固定官方源或传入已校验的本地源。完整工具环境与 29 组结果见\href{https://github.com/defineiocc02/20bit\_SAR\_ADC\_Behaviour\_Verification/blob/feat/ad4080-peripherals/subsys/ad4080\_peripherals/evidence/acceptance.md}{回归入口}。\par

已有 Verilator、C++ 编译器与安装了校准依赖的 Python 时，可重跑工程算术与图 96 判别，并保存机器可读结果。最后一条在原 DSP 仍给出 1024 时返回 2。\par

\begin{TraceCodeBlock}
\TraceCodeLine{p=\penalty150\relax{}subsys/\penalty150\relax{}ad4080\_\penalty150\relax{}peripherals}
\TraceCodeLine{python\hspace*{\fontdimen2\font}\allowbreak{}-\penalty150\relax{}m\hspace*{\fontdimen2\font}\allowbreak{}pip\hspace*{\fontdimen2\font}\allowbreak{}install\hspace*{\fontdimen2\font}\allowbreak{}-\penalty150\relax{}r\hspace*{\fontdimen2\font}\allowbreak{}"\$p/\penalty150\relax{}calibration/\penalty150\relax{}requirements\penalty150\relax{}.\penalty150\relax{}txt"}
\TraceCodeLine{python\hspace*{\fontdimen2\font}\allowbreak{}-\penalty150\relax{}B\hspace*{\fontdimen2\font}\allowbreak{}-\penalty150\relax{}m\hspace*{\fontdimen2\font}\allowbreak{}unittest\hspace*{\fontdimen2\font}\allowbreak{}discover\hspace*{\fontdimen2\font}\allowbreak{}-\penalty150\relax{}s\hspace*{\fontdimen2\font}\allowbreak{}"\$p/\penalty150\relax{}calibration"\penalty150\relax{}\hspace*{\fontdimen2\font}\allowbreak{}-\penalty150\relax{}p\hspace*{\fontdimen2\font}\allowbreak{}'test\_\penalty150\relax{}*.\penalty150\relax{}py'}
\TraceCodeLine{python\hspace*{\fontdimen2\font}\allowbreak{}-\penalty150\relax{}B\hspace*{\fontdimen2\font}\allowbreak{}"\$p/\penalty150\relax{}calibration/\penalty150\relax{}verify\_\penalty150\relax{}rtl.\penalty150\relax{}py"\hspace*{\fontdimen2\font}\allowbreak{}-\penalty150\relax{}-\penalty150\relax{}repo\hspace*{\fontdimen2\font}\allowbreak{}.\penalty150\relax{}\hspace*{\fontdimen2\font}\allowbreak{}-\penalty150\relax{}-\penalty150\relax{}out\hspace*{\fontdimen2\font}\allowbreak{}/\penalty150\relax{}tmp/\penalty150\relax{}ad4080-\penalty150\relax{}order.\penalty150\relax{}json\hspace*{\fontdimen2\font}\allowbreak{}-\penalty150\relax{}-\penalty150\relax{}require-\penalty150\relax{}manual-\penalty150\relax{}order}
\end{TraceCodeBlock}

复核手册版本时可自行从 ADI 下载并校验原件。量测根目录和临时文件均不进入生产 RTL。\par

\begin{TraceCodeBlock}
\TraceCodeLine{curl\hspace*{\fontdimen2\font}\allowbreak{}-\penalty150\relax{}fL\hspace*{\fontdimen2\font}\allowbreak{}-\penalty150\relax{}-\penalty150\relax{}retry\hspace*{\fontdimen2\font}\allowbreak{}3\hspace*{\fontdimen2\font}\allowbreak{}-\penalty150\relax{}o\hspace*{\fontdimen2\font}\allowbreak{}/\penalty150\relax{}tmp/\penalty150\relax{}ad4080-\penalty150\relax{}manual.\penalty150\relax{}pdf\hspace*{\fontdimen2\font}\allowbreak{}https:\penalty150\relax{}/\penalty150\relax{}/\penalty150\relax{}www.\penalty150\relax{}analog.\penalty150\relax{}com/\penalty150\relax{}media/\penalty150\relax{}en/\penalty150\relax{}technical-\penalty150\relax{}documentatio\penalty150\relax{}n/\penalty150\relax{}data-\penalty150\relax{}sheets/\penalty150\relax{}ad4080.\penalty150\relax{}pdf}
\TraceCodeLine{shasum\hspace*{\fontdimen2\font}\allowbreak{}-\penalty150\relax{}a\hspace*{\fontdimen2\font}\allowbreak{}256\hspace*{\fontdimen2\font}\allowbreak{}/\penalty150\relax{}tmp/\penalty150\relax{}ad4080-\penalty150\relax{}manual.\penalty150\relax{}pdf}
\end{TraceCodeBlock}
\end{document}
